% Part: lambda-calculus
% Chapter: church-rosser
% Section: parallel-beta-reduction

\documentclass[../../../include/open-logic-section]{subfiles}

\begin{document}

\olfileid{lam}{cr}{pb}

\olsection{\texorpdfstring{并行 $\beta$-归约}{并行 beta-归约}}

我们引入\emph{并行 $\beta$-归约}的概念，并证明其具有 Church--Rosser 性质。

\begin{defn}[并行 $\beta$-归约，$\bredpar$]\ollabel{defn:bredpar}
  项上的并行归约（$\bredpar$）归纳定义如下：
  \begin{enumerate}
    \item \ollabel{defn:bredpar1} $x \bredpar x$。
    \item \ollabel{defn:bredpar2} 若 $N \xrightarrow{\beta} N'$，则 $\lambd[x][N] \bredpar
      \lambd[x][N']$。
    \item \ollabel{defn:bredpar3} 若 $P \bredpar P'$ 且 $Q \bredpar Q'$，则 $PQ \bredpar P'Q'$。
    \item \ollabel{defn:bredpar4} 若 $N \bredpar N'$ 且 $Q \bredpar Q'$，则 $(\lambd[x][N])Q \bredpar \Subst{N'}{Q'}{x}$。
  \end{enumerate}
\end{defn}

并行 $\beta$-归约允许我们在一步中归约项内的任意多个可归约式。其与 $\beta$-归约的区别在于，我们只能收缩原始项中出现的可归约式，不能收缩因并行 $\beta$-归约而产生的可归约式。例如，项 $(\lambd[f][fx])(\lambd[y][y])$ 只能并行 $\beta$-归约到自身或 $(\lambd[y][y])x$，而不能进一步归约到~$x$，尽管它可 $\beta$-归约到~$x$，因为该可归约式仅在一步并行 $\beta$-归约之后才会产生。不过，第二步并行 $\beta$-归约即可得到~$x$。

\begin{thm}\ollabel{thm:refl}
  $M \bredpar M$。
\end{thm}

\begin{proof}
  练习。
\end{proof}

\begin{prob}
  证明 \olref[lam][cr][pb]{thm:refl}。
\end{prob}

\begin{defn}[$\beta$-完全发展]\ollabel{defn:bcd}
  项 $M$ 的\emph{$\beta$-完全发展}~$\bcd{M}$ 归纳定义如下：
  \begin{align}
    \bcd{x} &= x \ollabel{defn:bcd1} \\
    \bcd{(\lambd[x][N])} &= \lambd[x][\bcd{N}] \ollabel{defn:bcd2}\\
    \bcd{(PQ)} &= \bcd{P}\bcd{Q} && \text{若 $P$ 不是 $\lambd$-抽象}
    \ollabel{defn:bcd3} \\
    \bcd{((\lambd[x][N])Q)} &= \Subst{\bcd{N}}{\bcd{Q}}{x} \ollabel{defn:bcd4}
  \end{align}
\end{defn}

一个项的 $\beta$-完全发展，顾名思义，是一次“完全的并行归约”。对于并行 $\beta$-归约，我们仍可选择不收缩某个可归约式，但对于完全发展，我们只能收缩所有可归约式。因此，$(\lambd[f][fx])(\lambd[y][y])$ 的完全发展是 $(\lambd[y][y])x$，而不是它自身。

\begin{editorial}
  这个定义存在一个问题：我们尚未介绍如何在 ($\lambd$-)项上递归地定义函数。将在未来修正。
\end{editorial}

\begin{lem}\ollabel{lem:comp}
  若 $M \bredpar M'$ 且 $R \bredpar R'$，则 $\Subst{M}{R}{y}
  \bredpar \Subst{M'}{R'}{y}$。
\end{lem}

\begin{proof}
  对 $M \bredpar M'$ 的推导进行归纳。
  \begin{enumerate}
    \item 最后一步是 \olref{defn:bredpar1}：练习。
    \item 最后一步是 \olref{defn:bredpar2}：那么 $M$ 是 $\lambd[x][N]$ 且 $M'$ 是 $\lambd[x][N']$，其中 $N \bredpar N'$。要证 $\Subst{(\lambd[x][N])}{R}{y} \bredpar
      \Subst{(\lambd[x][N'])}{R'}{y}$，即 $\lambd[x][\Subst{N}{R}{y}] \bredpar
      \lambd[x][\Subst{N'}{R}{y}]$。这由 \olref{defn:bredpar2} 及归纳假设直接可得。
    \item 最后一步是 \olref{defn:bredpar3}：练习。
    \item 最后一步是 \olref{defn:bredpar4}：$M$ 是 $(\lambd[x][N])Q$ 且 $M'$ 是 $\Subst{N'}{Q'}{x}$。要证 $\Subst{((\lambd[x][N])Q)}{R}{y} \bredpar
      \Subst{\Subst{N'}{Q'}{x}}{R'}{y}$，即 $(\lambd[x][\Subst{N}{R}{y}])\Subst{Q}{R}{y} \bredpar
      \Subst{\Subst{N'}{R'}{y}}{\Subst{Q'}{R'}{y}}{x}$。这由 \olref{defn:bredpar4} 及归纳假设可得。
  \end{enumerate}
\end{proof}

\begin{prob}
  完成 \olref[lam][cr][pb]{lem:comp} 的证明。
\end{prob}

\begin{lem}\ollabel{lem:cont}
  若 $M \bredpar M'$，则 $M' \bredpar \bcd{M}$。
\end{lem}

\begin{proof}
  对 $M \bredpar M'$ 的推导进行归纳。
  \begin{enumerate}
    \item 最后一条规则是 \olref{defn:bredpar1}：练习。
    \item 最后一条规则是 \olref{defn:bredpar2}：$M$ 为
      $\lambd[x][N]$ 且 $M'$ 为 $\lambd[x][N']$，其中
      $N \bredpar N'$。要证 $\lambd[x][N'] \bredpar
      \bcd{(\lambd[x][N])}$，即 $\lambd[x][N'] \bredpar
      \lambd[x][\bcd{N}]$（据 \olref{defn:bcd2}）。这由
      \olref{defn:bredpar2} 及归纳假设可得。
    \item 最后一条规则是 \olref{defn:bredpar3}：$M$ 为 $PQ$ 且 $M'$
      为 $P'Q'$，其中 $P$, $Q$, $P'$, $Q'$ 满足 $P \bredpar P'$
      且 $Q \bredpar Q'$。由归纳假设，$P'
      \bredpar \bcd{P}$ 且 $Q' \bredpar \bcd{Q}$。
      \begin{enumerate}
        \item 若存在 $x$ 和 $N$ 使得 $P$ 为 $\lambd[x][N]$，则
          存在 $N'$ 使得 $P'$ 为 $\lambd[x][N']$ 且 $N \bredpar N'$。由归纳假设，$N' \bredpar \bcd{N}$ 且
          $Q' \bredpar \bcd{Q}$。于是由 \olref{defn:bredpar4} 得 $(\lambd[x][N'])Q' \bredpar
          \Subst{\bcd{N}}{\bcd{Q}}{x}$。
        \item 若 $P$ 不是 $\lambd$-抽象，则由 \olref{defn:bredpar3} 得 $P'Q' \bredpar
          \bcd{P}\bcd{Q}$，且根据 \olref{defn:bcd3}，其右侧即为 $\bcd{PQ}$。
      \end{enumerate}
    \item 最后一条规则是 \olref{defn:bredpar4}：$M$ 为
      $(\lambd[x][N])Q$ 且 $M'$ 为 $\Subst{N'}{Q'}{x}$，且 $N \bredpar N'$、$Q \bredpar
      Q'$。由归纳假设，$N' \bredpar \bcd{N}$ 且
      $Q' \bredpar \bcd{Q}$。由 \olref{lem:comp} 得
      $\Subst{N'}{Q'}{x} \bredpar \Subst{\bcd{N}}{\bcd{Q}}{x}$，
      其右侧恰好是 $\bcd{((\lambd[x][N])Q)}$。
  \end{enumerate}
\end{proof}

\begin{prob}
  完成 \olref[lam][cr][pb]{lem:cont} 的证明。
\end{prob}

\begin{thm}\ollabel{thm:cr}
  $\bredpar$ 具有 Church--Rosser 性质。
\end{thm}

\begin{proof}
  由 \olref{lem:cont} 直接可得。
\end{proof}

\end{document}
